% TOP_PROBLEM: {"id": 1295, "title": "Agoh-Giuga Conjecture", "category_id": 1, "category": {"id": 1, "name": "number_theory", "display_name": "Number Theory", "description": "Properties of integers, prime numbers, Diophantine equations.", "slug": "number-theory", "order_index": 1, "created_at": "2026-07-31T15:26:25.670Z"}, "status": "open"}
% ATTEMPT_STATUS: unresolved
% Research attempt by GPT_6_Astra_Ultra, 1 October 2026.
% Canonical UnsolvedMath ID; original statements and sources preserved.
% FIRST_POSTED: 2026-10-01
% Imported from: unsolvedmath_additional_500.tex; UnsolvedMath ID 1295
% !TeX program = lualatex
% Compile twice: lualatex 1295.tex
% Self-contained: no external bibliography, images, or \input files.
% Numeric section labels and visible numbers are original UnsolvedMath IDs.
\documentclass[11pt,a4paper]{article}
\usepackage[margin=24mm,headheight=14pt]{geometry}
\usepackage{fontspec}
\setmainfont{Latin Modern Roman}
\setsansfont{Latin Modern Sans}
\setmonofont{Latin Modern Mono}
\usepackage{amsmath,amssymb,mathtools}
\usepackage{unicode-math}
\setmathfont{Latin Modern Math}
\usepackage{microtype}
\usepackage{enumitem,xurl,needspace}
\usepackage[dvipsnames]{xcolor}
\usepackage{hyperref}
\hypersetup{unicode=true,colorlinks=true,linkcolor=MidnightBlue,urlcolor=MidnightBlue,
 pdftitle={UnsolvedMath Problem 1295},pdfauthor={GPT 6 Astra Ultra},bookmarksnumbered=true}
\usepackage{bookmark,tocloft,titlesec,fancyhdr}
\setlength{\cftsecnumwidth}{4.2em}
\cftsetpnumwidth{2.5em}
\cftsetrmarg{4em}
\titleformat{\section}{\Large\bfseries\raggedright}{\thesection}{0.7em}{}
\pagestyle{fancy}\fancyhf{}
\fancyhead[L]{\small\sffamily GPT 6 Astra Ultra research attempt}
\fancyhead[R]{\small Original catalogue IDs}
\fancyfoot[C]{\thepage}
\setlength{\parindent}{0pt}
\setlength{\parskip}{5pt plus 1pt}
\setlength{\emergencystretch}{3em}
\setcounter{tocdepth}{1}\setcounter{secnumdepth}{1}
\setlist[itemize]{leftmargin=3em,itemsep=4pt,topsep=4pt,parsep=0pt}
\allowdisplaybreaks
\newcommand{\N}{\mathbb{N}}
\newcommand{\Z}{\mathbb{Z}}
\newcommand{\Q}{\mathbb{Q}}
\newcommand{\R}{\mathbb{R}}
\newcommand{\C}{\mathbb{C}}
\newcommand{\F}{\mathbb{F}}
\newcommand{\A}{\mathbb{A}}
\newcommand{\PP}{\mathbb{P}}
\newcommand{\E}{\mathbb{E}}
\newcommand{\eps}{\varepsilon}
\newcommand{\dd}{\,\mathrm d}
\newcommand{\sref}[2]{\hyperlink{source:#1:#2}{[S#2]}}
\newif\ifshowcatalogue
\showcataloguetrue
\newcommand{\cataloguescope}[1]{\ifshowcatalogue
 \paragraph{Statement recorded in the supplied source.}
 #1\fi}
\begin{document}
% UnsolvedMath ID 1295; source catalogue code NT-087

\setcounter{section}{1294}

\section{The Agoh--Giuga Primality Criterion}\label{1295}

{\small\sffamily UnsolvedMath ID 1295 · Number Theory}

\subsection{Definitions and mathematical statement}

\paragraph{Shared notation.}Unless a section says otherwise, $\N=\{1,2,\ldots\}$,
$\Z$ is the ring of integers, $\Q,\R,\C$ are the rational, real, and complex
fields, $[N]=\{1,\ldots,\lfloor N\rfloor\}$, and $\log$ is the natural
logarithm. For real functions of a parameter tending to infinity,
$f\ll g$ and $f=O(g)$ mean $|f|\leq Cg$ eventually for some fixed $C$;
$f\gg g$ reverses this comparison; $f=o(g)$ means $f/g\to0$;
$f\sim g$ means $f/g\to1$; and $f\asymp g$ means both order bounds.
A subscript on $O$ or $\ll$ specifies allowed constant dependence.
The expression $n^{o(1)}$ in an upper bound means growth smaller than
$n^\epsilon$ for each fixed $\epsilon>0$. Local definitions govern any
exception, finite-field convention, or use of the same symbol for another
object. An asymptotic claim is not automatically an effective algorithm.



Define Bernoulli numbers by $t/(e^t-1)=\sum_{j\geq0}B_jt^j/j!$, so $B_1=-1/2$. For an integer $n\geq2$, interpret $nB_{n-1}\equiv-1\pmod n$ in the localization of $\Z$ obtained by inverting integers coprime to $n$: the rational difference $nB_{n-1}+1$ must belong to $n\Z_{(n)}$. This states a criterion for arbitrary integers, not only for an input already assumed prime. \sref{1295}{1} \sref{1295}{2}

\paragraph{Statement recorded in the supplied source.}
Is $p$ prime if and only if $pB_{p-1} \equiv -1 \pmod{p}$ for the Bernoulli number $B_{p-1}$? \sref{1295}{1}

\paragraph{Residual task reported by the archive.}

The embedded review (2026-08-17) records: Prove the criterion fails for every composite n, or find a composite Carmichael--Giuga number satisfying it. \sref{1295}{1}

\subsection{Short English statement}

Does this congruence involving a Bernoulli number characterize prime integers exactly?

\subsection{Research attempt}

\paragraph{Scope and theorem input.}
The variable is an arbitrary integer $n\ge2$, and the rational
congruence is interpreted in the localization specified above. The
von Staudt--Clausen theorem and the primary equivalence paper of Kellner
[E1] were inspected on 1 October 2026. We use the standard theorem
\[
 B_{2k}=A-\sum_{q-1\mid2k}\frac1q,\qquad A\in\mathbb Z,
\]
where the sum is over primes. The attempt derives the exact local
conditions from it, keeping denominator and prime-power issues visible.

\paragraph{The even case and the Bernoulli convention.}
The generating function gives $B_1=-1/2$ and $B_j=0$ for odd
$j>1$: after adding $t/2$, the function $t/(e^t-1)$ is even.
Thus $n=2$ passes because $2B_1=-1$. Every even $n>2$ fails because
$B_{n-1}=0$, making $nB_{n-1}+1=1$, which is not divisible by $n$
in the localization. This excludes every even composite without any
bound. From now on take $n>1$ odd, so $n-1$ is positive and even.

\paragraph{A local valuation proof of squarefreeness.}
Fix $p^a\Vert n$. In the von Staudt--Clausen expansion, all summands
with denominators different from $p$ are integral at $p$. If
$p-1\nmid n-1$, there is no $1/p$ term, so $nB_{n-1}+1\equiv1\pmod p$
and the criterion fails. If $p-1\mid n-1$, the only nonintegral
term is $-1/p$, with no cancellation possible against $p$-integral
terms. Therefore $v_p(B_{n-1})=-1$.

If $a\ge2$, it follows that $nB_{n-1}$ is divisible by $p$, again
making its sum with one nonzero modulo $p$. Hence a passing integer
must be squarefree and must satisfy $p-1\mid n-1$ for every $p\mid n$.
This conclusion cannot be obtained just by reducing a rational numerator
modulo $n$ before checking its denominator; the localization is part of
the argument.

\paragraph{The remaining local condition is exactly Giuga's.}
Now suppose $n$ is squarefree and every $p\mid n$ satisfies
$p-1\mid n-1$. Multiplying the Bernoulli expansion by $n$ and
reducing locally modulo $p$ leaves only the term $-n/p$:
\[
                         nB_{n-1}+1\equiv1-\frac np\pmod p.
\]
All other denominators are units at $p$. Thus the rational congruence
is equivalent to
\[
 p-1\mid n-1\quad\hbox{and}\quad p\mid n/p-1
                    \qquad\text{for every }p\mid n,
\]
together with squarefreeness. Conversely these conditions make the
rational difference integral at every prime dividing $n$ and divisible
by each such prime; because $n$ is squarefree, they give membership
in $n\mathbb Z_{(n)}$. This proves both directions for odd $n$.

For prime $n=p$, both conditions are immediate, proving the prime
direction of the conjecture. For composite $n$, the second is the
Giuga condition, while the first combined with squarefreeness is
Korselt's criterion for a Carmichael number. The simultaneous
requirement is much stronger than either condition separately.

\paragraph{What small-support arguments can exclude.}
The Giuga condition alone already excludes products of two distinct
primes: if $n=pq$ with $p<q$, it would require $q\mid p-1$.
For three distinct odd primes, its reciprocal form gives another
obstruction. Squarefreeness and the local divisibilities imply
\[
                   \sum_{p\mid n}\frac1p-\frac1n
                       \text{ is a positive integer}.
\]
But for three odd primes this is smaller than
$1/3+1/5+1/7<1$. Hence a composite passing the Bernoulli criterion
must have at least four prime factors by this elementary argument.
This modest lower bound is not claimed to be competitive with the
literature; it makes explicit that an unbounded-support case survives.

The number thirty illustrates why conventions cannot be mixed.
It satisfies the Giuga divisibilities, but it is even and fails the
Bernoulli criterion. An infinite family of Giuga numbers, if one were
proved, would not itself produce a counterexample unless at least one
also met every $p-1\mid n-1$ condition.

\paragraph{Executed rational check.}
The script constructs $B_0,\ldots,B_{99}$ as exact rational numbers
from
\[
 B_m=-\frac1{m+1}\sum_{j=0}^{m-1}\binom{m+1}{j}B_j.
\]
For each $2\le n\le100$, it reduces $nB_{n-1}+1$ to lowest terms,
requires its denominator to be coprime to $n$, and requires its numerator
to be divisible by $n$. Exactly the primes pass. Files are
\texttt{checks\_second.py}, \texttt{results\_second.json} under
\path{_Local/Erdos_problems/UnsolvedMath/attempt_audit_2026_10_01/number_theory/}.
The calculation checks the rational formulation itself, with its sign
convention, rather than only a related power-sum test.

\paragraph{Remaining gap and outcome.}
The attempt proves the prime direction and reduces every possible
counterexample to an odd squarefree integer satisfying both displayed
families of divisibilities. It does not show that these families are
globally incompatible for arbitrarily many prime factors.
\textbf{Outcome: UNRESOLVED.} The equivalence is a reduction of the
problem, not a proof that composite examples cannot occur.

\subsection{Sources}

{\small

\begin{itemize}

\item[\hypertarget{source:1295:1}{S1}] ULAM.AI, \emph{UnsolvedMath}, supplied \texttt{problems.json}, record 1295 (NT-087), statement and background. The embedded literature-review date is 2026-08-17; that is the archive’s review date, not a fresh certification. Dataset landing page: \url{https://huggingface.co/datasets/ulamai/UnsolvedMath}.

\item[\hypertarget{source:1295:2}{S2}] D. Borwein, J. M. Borwein, P. B. Borwein and R. Girgensohn, Giuga's conjecture on primality, American Mathematical Monthly 103 (1996), 40--50. \url{https://doi.org/10.2307/2975190}. Bibliographic information imported from the supplied record.

\item[\hypertarget{source:1295:3}{S3}] Agoh--Giuga conjecture overview. \url{https://en.wikipedia.org/wiki/Agoh%E2%80%93Giuga_conjecture}. Bibliographic information imported from the supplied record.


\item[E1] Bernd C. Kellner, \emph{The Equivalence of Giuga's and
Agoh's Conjectures}, arXiv:math/0409259.
\url{https://arxiv.org/pdf/math/0409259}.
Primary paper inspected 1 October 2026. The von Staudt--Clausen theorem
is a stated input; the local valuation deductions are spelled out above.

\end{itemize}

}

\label{end:1295}
\end{document}
